\chapter*{Conclusion Générale}
\addcontentsline{toc}{chapter}{Conclusion Générale}

Le projet de fin d'année présenté tout au long de ce rapport s'est articulé autour de la conception et du développement d'une application web dédiée à la gestion des immobilisations pour la Société Tunisienne de l'Électricité et du Gaz. L'objectif initial consistait à répondre à une problématique concrète de suivi et de centralisation du patrimoine matériel, informatique et roulant, afin de pallier les limites des méthodes de gestion traditionnelles.

Les différentes étapes menées à bien ont permis d'aboutir à une solution logicielle fonctionnelle et sécurisée. La phase d'analyse des exigences a posé les bases indispensables en identifiant avec précision les besoins des différents services et profils d'utilisateurs. L'étape de conception, appuyée sur une modélisation rigoureuse des structures et des processus, a garanti l'élaboration d'une architecture de base de données relationnelle stable et cohérente. Enfin, la phase d'implémentation a concrétisé ces travaux à travers une interface web intuitive, interconnectée à une logique applicative robuste capable d'assurer la traçabilité complète de chaque actif au sein des locaux.

Sur le plan personnel et technique, ce travail a constitué une opportunité d'approfondir et de consolider des compétences pratiques en développement complet et en gestion de bases de données, tout en se confrontant aux réalités opérationnelles et aux exigences de sécurité d'une grande entreprise publique nationale. L'immersion au sein du Département d'Exploitation Informatique Sud a permis de comprendre l'importance d'une planification structurée et d'une conception modulaire pour garantir la réussite d'un projet informatique.

Bien que l'application réponde aux spécifications établies dans le cahier des charges, le système reste ouvert à de futures extensions opérationnelles. En guise de perspectives, l'intégration d'un module de notification automatique pourrait être envisagée pour planifier les opérations de maintenance préventive du matériel électrique et le contrôle technique des véhicules. De plus, l'adoption de technologies d'identification rapide, telles que la génération et la lecture de codes-barres ou de codes QR pour chaque immobilisation, permettrait d'automatiser entièrement le processus de contrôle physique lors des campagnes d'inventaire annuel.